\section{Interpretability as Hypothesis Testing}
\label{sec:interp-hypothesis}

Mechanistic interpretability is currently largely exploratory. The BP
framework changes this: it generates specific, testable predictions about
what trained transformers should look like at the weight level.

\subsection*{Specific Predictions}

\begin{itemize}
\item Boolean-AND heads should have approximately rank-1 Q/K circuits
consistent with \texttt{projectDim}-style matching.
\item The dominant singular vector of the Q/K circuit should correspond to
a monosemantic SAE feature.
\item The V circuit should show \texttt{crossProject}-style routing from a
belief-encoding dimension to a scratch-slot dimension.
\item $W_1$ rows of the FFN should cluster by input SAE feature.
\item $W_2$ columns should correspond to interpretable belief updates.
\item The relationship between input and output SAE features for each FFN
unit should be interpretable as a conditional probability table entry.
\end{itemize}

These are specific structural claims that can be confirmed or falsified by
inspecting the weights of a trained model with a good SAE.

\subsection*{What a Negative Result Would Mean}

If boolean-AND heads do not show rank-1 Q/K structure, that would falsify
the constructive BP account for those heads --- not just complicate it. The
BP framework is a genuine scientific hypothesis because it makes predictions
that could be wrong.

\subsection*{The Closing Experiment}

Take a model with a trained SAE. Identify boolean-AND heads by behavior.
Project the Q/K matrix onto the SAE feature basis. Check if the dominant
component corresponds to a monosemantic feature. Check if the V circuit
routes from a belief dimension to a scratch slot. Check if input/output SAE
feature pairs for active FFN units are interpretable as conditional
relationships. If this succeeds, the implicit factor graph is legible.

\subsection*{Epistemic Status}

\begin{center}
\begin{tabular}{ll}
\toprule
Claim & Status \\
\midrule
BP predicts rank-1 Q/K structure in AND heads & Empirical hypothesis \\
SAE features correspond to factor graph nodes & Empirical hypothesis \\
Factor graph is legible from weights + SAE & Empirical hypothesis \\
Negative SAE result would falsify BP account & Direct consequence \\
\bottomrule
\end{tabular}
\end{center}